\section{Proofs}
\label{app:proofs}

This appendix proves the five propositions of the main text and states, for each,
exactly where the proof stops. Two of them are unconditional
(Propositions~\ref{prop:disjoint} and~\ref{prop:cheb}); two are conditional on a
margin-density constant $\kappa$ that we measure rather than assume
(Propositions~\ref{prop:cheb-acc} and~\ref{prop:amp}); one is proved in the
special case $d=1$ that covers every scorable multi-task scope on our benchmark
and is open beyond it (Proposition~\ref{prop:qm}). We also record, in
Appendix~\ref{app:narrowed}, three claims that earlier drafts of this work stated
too strongly, together with the counterexamples that forced the narrowing. They
are kept visible because each one changes how a number in the main text may be
reported.

\subsection{Setting and gauge}
\label{app:gauge}

Fix a scope $S$ with label set $Y_S=\{\ell_1,\dots,\ell_{K}\}$, $K=K_S$. For an
input $\bx$ and adapter parameters $\theta$, write the verbalizer-restricted logit
vector
\begin{equation}
  \bm{\lambda}_\theta(\bx)\in\R^{K},
  \qquad
  \big[\bm{\lambda}_\theta(\bx)\big]_k = \bu_{\ell_k}^\top\bh_\theta(\bx),
  \label{eq:app-logit}
\end{equation}
with $\bu_y$ the frozen unembedding row for token $y$. An offset $\bb\in\R^{K}$
acts by $\hat y(\bx)=\arg\max_k\big([\bm{\lambda}_\theta(\bx)]_k+b_k\big)$.

\paragraph{The gauge.}
For any $c\in\R$, the offsets $\bb$ and $\bb+c\mathbf{1}$ induce the same
restricted argmax at every input, so $\bb$ acts only through its image in the
quotient $\R^{K}/\mathrm{span}(\mathbf{1})$. We fix the gauge by the zero-sum
projector
\begin{equation}
  P_0\bb \;=\; \bb - \Big(\tfrac{1}{K}\textstyle\sum_k b_k\Big)\mathbf{1},
  \qquad
  V_0 \;=\; \{\bm{v}\in\R^{K}:\ \textstyle\sum_k v_k = 0\}\ \cong\ \R^{K-1},
  \label{eq:app-gauge}
\end{equation}
and equip $V_0$ with the Euclidean norm. Every offset below is an element of
$V_0$. This is the same $P_0$ as in Section~\ref{sec:method}, and it is the reason
Equation~\eqref{eq:cost} counts $K_S-1$ rather than $K_S$ scalars per stored
vector.

For $K=2$ the space $V_0$ is one-dimensional: $P_0\bb=(-t,t)$ with
$t=(b_1-b_0)/2$, so an offset is fully described by the scalar contrast
$s=b_1-b_0$, and $\|P_0\bb-P_0\bb'\|_2=|s-s'|/\sqrt2$. This constant $1/\sqrt2$ is
a unit convention, not a dimension: it is why a scope whose four member contrasts
span a range of exactly $2$ is recorded with $q_1=0.7071$.

Write $R_g(\bb)$ for task $g$'s balanced accuracy under shared offset $\bb$ at the
current parameters, and
\begin{equation}
  \bz_g \in \argmax_{\bm{v}\in V_0} R_g(\bm{v}),
  \qquad
  R_g^{\mathrm{orc}} = R_g(\bz_g).
  \label{eq:app-orc}
\end{equation}
Ties in \eqref{eq:app-orc} are broken by any fixed rule; nothing below depends on
which, because the statements quantify over all shared $\bb$ and refer to $\bz_g$
only through its gauge-fixed value.

\subsection{Proof of Proposition~\ref{prop:disjoint} (disjoint verbalizers give a zero gap)}
\label{app:proof-disjoint}

\begin{proof}
Suppose the verbalizer sets $Y_1,\dots,Y_T$ are pairwise disjoint. An offset in
the sense of Section~\ref{sec:setup} is a vector $\bb\in\R^{V}$ over the whole
vocabulary, and task $g$'s prediction reads only the coordinates $\{b_y:y\in Y_g\}$.
Let $\bb_g^\star\in\R^{V}$ be an optimal offset for task $g$, supported on $Y_g$
(any optimal offset may be restricted to $Y_g$ without changing $g$'s predictions,
since coordinates outside $Y_g$ never enter $g$'s argmax). Define
\begin{equation}
  \bb^\star \;=\; \sum_{g=1}^{T}\bb_g^\star .
  \label{eq:app-disjoint-sum}
\end{equation}
Because the supports are pairwise disjoint, for every $g$ and every $y\in Y_g$ we
have $b^\star_y = [\bb_g^\star]_y$. Hence the restricted argmax of task $g$ under
$\bb^\star$ agrees with its restricted argmax under $\bb_g^\star$ at every input,
so $R_g(\bb^\star)=R_g^{\mathrm{orc}}$ for all $g$ simultaneously. Since
$R_g^{\mathrm{shr}}$ is the accuracy under the \emph{best} single shared offset, it
is at least $R_g(\bb^\star)$, and it is at most $R_g^{\mathrm{orc}}$ always. Both
inequalities together give $R_g^{\mathrm{shr}}=R_g^{\mathrm{orc}}$ and therefore
$\Delta_g^{\mathrm{id}}=0$ for every $g$ and every time $t$.
\end{proof}

\begin{remark}[The singleton corollary we use as a wiring check]
\label{rem:singleton}
The same argument applied to a single scope gives the form used in
Section~\ref{sec:method}: if $|\mathrm{mem}(S)|=1$ then $Z_S$ is a one-point set,
its Chebyshev centre \emph{is} that point, and $r_S=0$ exactly. This is a
prediction of the theory with no free constant, so any nonzero measured
$\Delta^{\mathrm{id}}$ on a singleton scope is a bug rather than a finding, and we
report the check in every run ($0$ violations, $87$--$144$ observations per arm,
worst absolute deviation $0.0$). Note what the check does \emph{not} cover: it
constrains the scoped rung, not a stored-offset method, because a stored optimum
is fitted at an earlier parameter value. That distinction is exactly what makes
the singleton split of Section~\ref{sec:sso} an isolate for staleness.
\end{remark}

\begin{remark}[Why this is the honest boundary, not a weakness]
Proposition~\ref{prop:disjoint} says the phenomenon we study has a precondition
that is a property of the benchmark and is checkable without running anything: the
tasks must reuse output tokens. On a stream with pairwise disjoint verbalizers
$\Delta^{\mathrm{id}}\equiv0$ and nothing in this paper applies. We state it first
for that reason.
\end{remark}

\subsection{Proof of Proposition~\ref{prop:cheb} (logit-space minimax)}
\label{app:proof-cheb}

Recall $Z_S=\{P_0\bz_g : g\in\mathrm{mem}(S)\}\subset V_0$, a finite set of
$n=|\mathrm{mem}(S)|$ points, and
$F(\bb)=\max_{g\in\mathrm{mem}(S)}\|\bb-P_0\bz_g\|_2$.

\begin{proof}
\emph{(i) The minimum over $\R^{K}$ is attained inside $V_0$.}
For $\bb\in\R^{K}$ decompose $\bb = P_0\bb + \beta\mathbf{1}$ with
$\beta=\frac1K\sum_k b_k$. Since every $\bm{p}\in Z_S$ lies in $V_0$ and
$\mathbf{1}\perp V_0$, Pythagoras gives
$\|\bb-\bm{p}\|_2^2 = \|P_0\bb-\bm{p}\|_2^2 + K\beta^2 \ge \|P_0\bb-\bm{p}\|_2^2$,
with equality iff $\beta=0$. So $F(\bb)\ge F(P_0\bb)$ with equality iff
$\bb\in V_0$, and it suffices to minimise $F$ over $V_0$.

\emph{(ii) The minimum is attained.}
$F$ is a finite maximum of convex functions, hence convex, and continuous. It is
coercive: $F(\bb)\ge\|\bb\|_2-\max_g\|P_0\bz_g\|_2\to\infty$. A continuous coercive
function on the finite-dimensional space $V_0$ attains its infimum, and by
definition that infimum is the Chebyshev radius $r_S$ of $Z_S$ and its minimiser
set is the set of Chebyshev centres. This gives \eqref{eq:cheb}.

\emph{(iii) The minimiser is unique.}
Suppose $\bb_1\neq\bb_2$ both attain $F=r_S$, and let
$\bar\bb=\tfrac12(\bb_1+\bb_2)$. For each $\bm{p}\in Z_S$, the parallelogram law
applied to $\tfrac12(\bb_1-\bm{p})+\tfrac12(\bb_2-\bm{p})$ gives
\begin{equation}
  \|\bar\bb-\bm{p}\|_2^2
  \;=\;\tfrac12\|\bb_1-\bm{p}\|_2^2+\tfrac12\|\bb_2-\bm{p}\|_2^2
      -\tfrac14\|\bb_1-\bb_2\|_2^2
  \;\le\; r_S^2-\tfrac14\|\bb_1-\bb_2\|_2^2
  \;<\;r_S^2 .
  \label{eq:app-parallelogram}
\end{equation}
Taking the maximum over $\bm{p}\in Z_S$ yields $F(\bar\bb)<r_S$, contradicting
minimality. Hence the Chebyshev centre is unique. Combined with (i), the minimiser
over all of $\R^{K}$ is unique as well, and it is the zero-mean representative.

\emph{(iv) The diameter bound.}
For any $g,g'\in\mathrm{mem}(S)$ and any $\bb\in V_0$, the triangle inequality gives
$\|P_0\bz_g-P_0\bz_{g'}\|_2\le\|\bb-P_0\bz_g\|_2+\|\bb-P_0\bz_{g'}\|_2\le 2F(\bb)$.
Maximising the left side over the pair and minimising the right over $\bb$ gives
$r_S\ge\tfrac12\mathrm{diam}(Z_S)$. In particular $r_S=0$ iff
$\mathrm{diam}(Z_S)=0$, i.e.\ iff all member optima coincide after gauge fixing.

\emph{(v) $r_S=q_1(S)$.}
The $m$-centre quantization radius of Section~\ref{sec:ladder} is
$q_m(Z_S)=\min_{|C|\le m}\max_{\bm{p}\in Z_S}\mathrm{dist}(\bm{p},C)$. At $m=1$ a
codebook is a single point $C=\{\bb\}$ and $\mathrm{dist}(\bm{p},C)=\|\bb-\bm{p}\|_2$,
so $q_1(Z_S)=\min_{\bb} F(\bb)=r_S$.
\end{proof}

\begin{remark}[Uniqueness is a statement about vectors, not about decision rules]
Part (iii) gives a unique minimising \emph{vector} in $V_0$. As an inference rule
the minimiser is unique only modulo $\mathrm{span}(\mathbf{1})$, by the gauge
freedom of Appendix~\ref{app:gauge}; the zero-mean representative is the
norm-minimal member of that equivalence class. This is why the method stores
$P_0\bz_g$ and not $\bz_g$, and why the cost \eqref{eq:cost} is $K_S-1$ per vector.
\end{remark}

\begin{remark}[What makes the target streamable]
\label{rem:streamable}
The right-hand side of \eqref{eq:cheb} is a function of the member optima $Z_S$
alone. It does not involve the members' examples, their logits, their sample sizes,
or their arrival order. That is the entire licence for the method of
Section~\ref{sec:method}: if the target depended on the data through anything more
than $Z_S$, storing a $K_S$-vector per task could not in principle suffice.
Proposition~\ref{prop:cheb} therefore establishes that the SSO \emph{target} is
representable at the stated cost --- and nothing more. It does not establish that
the stored vectors remain valid as $\theta$ moves, which is the empirical question
Section~\ref{sec:sso} answers in the negative.
\end{remark}

\subsection{Proposition~\ref{prop:cheb-acc}: from distance to accuracy, and what it costs}
\label{app:proof-cheb-acc}

Proposition~\ref{prop:cheb} bounds a \emph{distance}. Turning it into a statement
about accuracy requires an assumption, and the reason is not technical: a distance
can be large while costing nothing. If every example of every member task sits far
from its decision boundary, moving the shared offset by a bounded amount flips no
prediction and each $R_g$ is locally constant. We therefore separate two
quantities that an earlier draft of this work conflated.

\begin{definition}[Disagreement and accuracy-decay profiles]
\label{def:profiles}
Fix a task $g$ and let $\hat y(\cdot\,;\bb)$ denote its restricted argmax under
offset $\bb$. For $t\ge0$ define
\begin{align}
  \Phi_g(t) &= \min_{\bm{u}\in V_0,\ \|\bm{u}\|_2=1}
    \ \Pr_{\bx\sim\cD_g}\!\big[\hat y(\bx;\bz_g+t\bm{u})\neq \hat y(\bx;\bz_g)\big],
    \label{eq:app-phi}\\
  \Psi_g(t) &= \min_{\bm{u}\in V_0,\ \|\bm{u}\|_2=1}
    \ \big[R_g(\bz_g)-R_g(\bz_g+t\bm{u})\big].
    \label{eq:app-psi}
\end{align}
$\Phi_g$ is the smallest fraction of predictions that \emph{change} when the offset
is displaced a distance $t$ from $g$'s own optimum; $\Psi_g$ is the smallest
resulting \emph{loss of accuracy}. Both minima are over directions, so both are
worst-case over the direction the shared offset happens to lie in. By optimality
of $\bz_g$ in \eqref{eq:app-orc}, $\Psi_g(t)\ge0$ for all $t$.
\end{definition}

\begin{assumption}[Margin density, disagreement form]
\label{asm:kappa-dis}
There exist $\kappa_{\mathrm{dis}}>0$ and $t_0>0$ with
$\Phi_g(t)\ \ge\ \kappa_{\mathrm{dis}}\cdot\min(t,t_0)$ for every $t\ge0$ and every
$g\in\mathrm{mem}(S)$.
\end{assumption}

\begin{assumption}[Margin density, accuracy form]
\label{asm:kappa-acc}
There exist $\kappa>0$ and $t_0>0$ with
$\Psi_g(t)\ \ge\ \kappa\cdot\min(t,t_0)$ for every $t\ge0$ and every
$g\in\mathrm{mem}(S)$.
\end{assumption}

Assumption~\ref{asm:kappa-acc} is the hypothesis of
Proposition~\ref{prop:cheb-acc} as stated in Section~\ref{sec:theory}; the
interval $[0,\varepsilon]$ there is $[0,t_0]$ here. The proof is then two lines,
and this is deliberate: the work is done by \emph{where} the two statements live,
not by the inequality itself.

\begin{proof}[Proof of Proposition~\ref{prop:cheb-acc}]
Let $\bb\in V_0$ be any shared offset on $S$. By Proposition~\ref{prop:cheb} there
is a member $g$ with $t:=\|\bb-P_0\bz_g\|_2\ \ge\ r_S$. Set
$\bm{u}=(\bb-P_0\bz_g)/t$ if $t>0$ (if $t=0$ then $r_S=0$ and the claim is
vacuous). Since $\Psi_g$ minimises over unit directions,
\begin{equation}
  R_g^{\mathrm{orc}}-R_g(\bb)
  \;=\; R_g(\bz_g)-R_g(\bz_g+t\bm{u})
  \;\ge\; \Psi_g(t)
  \;\ge\; \kappa\min(t,t_0)
  \;\ge\; \kappa\min(r_S,t_0),
  \label{eq:app-cheb-acc-chain}
\end{equation}
using Assumption~\ref{asm:kappa-acc} and then $t\ge r_S$ together with
monotonicity of $\tau\mapsto\min(\tau,t_0)$. Taking the maximum over
$g\in\mathrm{mem}(S)$ on the left can only increase it, which is
\eqref{eq:cheb-acc}.
\end{proof}

\noindent The same argument with $\Phi$ in place of $\Psi$ gives the disagreement
form, which we record separately because it is the one our measurement certifies.

\begin{lemma}[Disagreement-space minimax]
\label{lem:dis-minimax}
Under Assumption~\ref{asm:kappa-dis}, for every shared $\bb$ on $S$,
\begin{equation}
  \max_{g\in\mathrm{mem}(S)}
    \Pr_{\bx\sim\cD_g}\!\big[\hat y(\bx;\bb)\neq\hat y(\bx;\bz_g)\big]
  \;\ge\; \kappa_{\mathrm{dis}}\cdot\min(r_S,t_0).
\end{equation}
\end{lemma}

\subsubsection{What the measured constant certifies, stated precisely}
\label{app:kappa-measured}

We report $\kappa\approx0.0182$ (worst task, WiC) in Section~\ref{sec:theory}. The
estimator behind that number is a top-$2$ margin CDF, so it estimates
$\kappa_{\mathrm{dis}}$ of Assumption~\ref{asm:kappa-dis}, not $\kappa$ of
Assumption~\ref{asm:kappa-acc}. The two are ordered, and only in one direction:

\begin{lemma}[The disagreement constant bounds the accuracy constant]
\label{lem:kappa-order}
$\Psi_g(t)\le\Phi_g(t)$ for all $t\ge0$. Consequently the largest valid constant in
Assumption~\ref{asm:kappa-acc} is at most the largest valid constant in
Assumption~\ref{asm:kappa-dis}: $\kappa\le\kappa_{\mathrm{dis}}$.
\end{lemma}

\begin{proof}
Fix $g$ and $t$, and for a unit direction $\bm{u}$ write
$\Psi_g^{\bm{u}}(t)=R_g(\bz_g)-R_g(\bz_g+t\bm{u})$ and
$D_g^{\bm{u}}(t)=\Pr[\hat y(\bx;\bz_g+t\bm{u})\neq\hat y(\bx;\bz_g)]$.

\emph{Step 1 (pointwise in $\bm{u}$).} Let $A$ be the event that the prediction is
correct at $\bz_g$ and incorrect at $\bz_g+t\bm{u}$, and $B$ the reverse. Then
$\Psi_g^{\bm{u}}(t)=\Pr[A]-\Pr[B]\le\Pr[A]\le\Pr[A\cup B]\le D_g^{\bm{u}}(t)$,
since $A$ and $B$ are disjoint subsets of the disagreement event.

\emph{Step 2 (transfer to the profiles).} Let $\bm{u}^\star$ attain
$\Phi_g(t)=\min_{\bm{u}}D_g^{\bm{u}}(t)$, which exists because the unit sphere in
$V_0$ is compact and $D_g^{\bm{u}}$ takes finitely many values on a finite sample.
Then
\begin{equation}
  \Psi_g(t)\;=\;\min_{\bm{u}}\Psi_g^{\bm{u}}(t)
  \;\le\;\Psi_g^{\bm{u}^\star}(t)
  \;\le\;D_g^{\bm{u}^\star}(t)
  \;=\;\Phi_g(t).
\end{equation}

\emph{Step 3.} If Assumption~\ref{asm:kappa-acc} holds with constant $\kappa$ then
$\kappa\min(t,t_0)\le\Psi_g(t)\le\Phi_g(t)$, so Assumption~\ref{asm:kappa-dis}
holds with the same constant. Hence the supremum of valid $\kappa$ is no larger
than the supremum of valid $\kappa_{\mathrm{dis}}$.
\end{proof}

\noindent So the number we measured is an \emph{upper} estimate of the constant that
Proposition~\ref{prop:cheb-acc} needs. We state the consequence rather than leaving
it implicit: \textbf{Equation~\eqref{eq:cheb-acc} evaluated at
$\kappa=0.0182$ is an optimistic reading of a bound we have not separately
certified, and Lemma~\ref{lem:dis-minimax} at
$\kappa_{\mathrm{dis}}=0.0182$ is the statement our measurement supports.} Nothing
downstream changes, because no empirical claim in this paper is derived from
\eqref{eq:cheb-acc}: the geometric half carries the load and every accuracy-space
number we report is measured. But the distinction has to be on the page, because
$0.0182$ appears next to an accuracy-space inequality in the main text.

Three further gaps between the assumption and the estimator, in the order they
matter:

\begin{enumerate}
  \item \textbf{Centring.} The estimator computes the margin CDF of the raw
    verbalizer logits, i.e.\ the flip curve around $\bb=\mathbf{0}$, whereas
    Definition~\ref{def:profiles} centres at $\bz_g$. The two differ by a shift of
    the margin distribution and \emph{we do not know the sign of the discrepancy};
    it is a proxy, not a conservative one.
  \item \textbf{Grid safety.} The curve is recorded on
    $t\in\{0.1,0.25,0.5,1,1.5,2,3,4\}$. The naive slope $\min_i N(t_i)/t_i$
    verifies $N(t)\ge\kappa t$ only at grid points, while the assumption needs all
    $t$; since $N$ is non-decreasing the worst case in $[t_i,t_{i+1})$ is
    $t\to t_{i+1}^-$ with $N$ still at $N(t_i)$, so the valid slope is
    $\min_i N(t_i)/t_{i+1}$. Grid ratios reach $2.5\times$ and the two columns
    differ by up to that factor. Only the latter is a lower bound, and it is what
    we report.
  \item \textbf{Multi-class conservatism.} For $K_S\ge3$ only top-$2$ flips are
    counted, so $N$ is itself a lower bound on the flip fraction and
    $\kappa_{\mathrm{dis}}$ is understated for MNLI, AGNews, DBpedia, Yahoo, Yelp
    and Amazon. This gap favours the bound.
\end{enumerate}

\noindent Item 2 is why the reported worst-task value is $0.0182$ rather than the
naive $0.0312$ or the older two-quantile estimate $0.0871$; the last should not be
cited. We quote the worst task rather than the median $0.0625$ because $\kappa$ must
hold \emph{simultaneously} for every member of a scope, so the smallest-$\kappa$
member sets the bound --- a median would overstate the usable constant by
$3.4\times$.

\subsubsection{The assumption is necessary, not decorative}
\label{app:kappa-necessary}

\begin{proposition}[No assumption-free accuracy-space form exists]
\label{prop:no-free-lunch}
There is no function $\omega$ with $\omega(\rho)>0$ for $\rho>0$ such that
$\max_{g\in\mathrm{mem}(S)}[R_g^{\mathrm{orc}}-R_g(\bb)]\ge\omega(r_S)$ holds for
every scope, every model and every shared $\bb$.
\end{proposition}

\begin{proof}
Two tasks share the binary verbalizer $\{\ell_0,\ell_1\}$, so $V_0\cong\R$ and an
offset is the scalar contrast $s$. Let task $1$ have every example at logit gap
$+c$ in favour of its true label and task $2$ every example at gap $-c$, for a
constant $c>0$. Then $z_1=0$ and any $z_2>c$ is optimal for task $2$, so
$r_S\ge c/2$. At $\bb=0$ task $1$ is perfect and task $2$ scores $0$, so the
left-hand side equals $1$ --- but it equals $1$ for \emph{every} $c$, including
$c\to0$, where $r_S\to0$. Conversely rescaling all logits by $\alpha$ multiplies
$r_S$ by $\alpha$ while leaving every argmax, hence every accuracy, unchanged. So
the left-hand side is constant along a family on which $r_S$ ranges over
$(0,\infty)$, and no positive $\omega(r_S)$ can bound it below for all members of
the family; taking $\alpha$ large also shows no upper bound of the form
$\omega(r_S)\le$ accuracy loss can be non-vacuous, since accuracy differences are
at most $1$ and $r_S$ is unbounded.
\end{proof}

\noindent This is the concrete form of the correction recorded in
Section~\ref{sec:theory}: an earlier version of our own analysis bounded an
accuracy difference below by a diameter-type quantity. Accuracy differences are at
most $1$; diameters are unbounded; the two are not comparable without a constant
that ties logit distance to flip probability. That constant is $\kappa$, and
Proposition~\ref{prop:no-free-lunch} is why it cannot be removed.

\begin{remark}[What the conditional proposition is still worth]
Given how close Assumption~\ref{asm:kappa-acc} is to the conclusion of
Proposition~\ref{prop:cheb-acc}, it is fair to ask what the proposition adds. The
answer is a change of quantifier, not a strengthening of the rate. The assumption
is \emph{local}, \emph{per-task}, and about a single task's own neighbourhood; the
conclusion is about the \emph{worst member of a scope} and its right-hand side
$r_S$ is a function of the member optima alone --- computable, by
Proposition~\ref{prop:cheb}, from the geometry of the demands without reference to
any model or dataset. That transfer is the content. We do not claim more for it,
and at $\kappa\approx0.018$ we do not use it quantitatively.
\end{remark}

\subsection{Proof of Proposition~\ref{prop:qm} (budgeted offsets)}
\label{app:proof-qm}

A budget of $m$ task-free offsets on a scope is a \emph{codebook}
$C=\{\bm{c}_1,\dots,\bm{c}_m\}\subset V_0$ together with a rule assigning each task
to a centre. Write $X=Z_S\subset V_0$ for the member optima, $d=\dim V_0=K_S-1$,
$T=|X|$, and
\begin{equation}
  q_m(X)\;=\;\min_{|C|\le m}\ \max_{\bm{x}\in X}\ \mathrm{dist}(\bm{x},C),
  \qquad \mathrm{dist}(\bm{x},C)=\min_{j}\|\bm{x}-\bm{c}_j\|_2 .
  \label{eq:app-qm}
\end{equation}
By construction $q_1\ge q_2\ge\cdots\ge q_T=0$, which is the monotonicity of the
ladder in Section~\ref{sec:ladder}. Note that \eqref{eq:app-qm} presumes
nearest-centre assignment; a task-free router that assigns otherwise can only do
worse, which is why $q_m$ is a bound on any routed method and not a promise about
one.

\subsubsection{The closed form at \texorpdfstring{$d=1$}{d=1}}

\begin{proposition}[Closed form, $m=2$, $d=1$]
\label{prop:q2-closed}
Let $X=\{x_1\le\cdots\le x_T\}\subset\R$ with $T\ge2$. Then
\begin{equation}
  q_2(X)\;=\;\min_{1\le j\le T-1}\ \max\!\Big(\tfrac{x_j-x_1}{2},\ \tfrac{x_T-x_{j+1}}{2}\Big),
  \label{eq:app-q2}
\end{equation}
attained by $C=\{\tfrac{x_1+x_j}{2},\ \tfrac{x_{j+1}+x_T}{2}\}$ for a minimising $j$.
\end{proposition}

\begin{proof}
\emph{Step 1: an optimal $2$-codebook induces a contiguous split.} Let
$C=\{c_1\le c_2\}$ and assign each point to its nearest centre. In $\R$ the
nearest-centre rule assigns $x$ to $c_1$ iff $x\le(c_1+c_2)/2$, so the assignment
is a threshold rule: there is $j\in\{0,\dots,T\}$ with $\{x_1,\dots,x_j\}\to c_1$
and $\{x_{j+1},\dots,x_T\}\to c_2$. Only the $T+1$ contiguous splits need be
considered. This is exactly the step that fails for $d\ge2$: there the Voronoi
cells of two centres are half-spaces, but enumerating orderings of the points does
not enumerate the achievable bipartitions.

\emph{Step 2: within a block the best centre is the midrange.} For a finite
$B\subset\R$, $\min_c\max_{x\in B}|x-c|=(\max B-\min B)/2$, attained uniquely at
$c=(\max B+\min B)/2$: for any $c$,
$\max_{x\in B}|x-c|\ge\tfrac12[(\max B-c)+(c-\min B)]=(\max B-\min B)/2$, with
equality iff $c$ is the midpoint.

\emph{Step 3: combine.} For split $j$ the covering radius is
$\max\big((x_j-x_1)/2,(x_T-x_{j+1})/2\big)$ by Step 2 applied to each block.
Minimising over $j$ gives \eqref{eq:app-q2}. The degenerate splits $j=0$ and $j=T$
leave one block empty and reduce to $q_1=(x_T-x_1)/2$, which is never better than
the best proper split, so restricting to $1\le j\le T-1$ is without loss.
\end{proof}

\begin{remark}[The equality clause, stated correctly]
\label{rem:equality-clause}
$q_2=q_1$ iff $T=1$ (both are $0$). For $T=2$ and $x_1\neq x_2$ we have
$q_2=0<q_1=(x_2-x_1)/2$: two centres cover two points exactly. An earlier draft of
this note wrote ``equality iff $T\le2$'', which is wrong, and wrong in a way that
would have contradicted our own measurement --- $T=2$ is precisely the
$\{\texttt{bad},\texttt{good}\}$ scope where the data records $q_2=0$. Left visible
rather than silently corrected.
\end{remark}

\subsubsection{The pigeonhole bound, and its tightness at \texorpdfstring{$d=1$}{d=1}}

A closed form requires solving the problem. For the paper we also want a bound that
reads off the geometry of the demands directly. Define the \emph{$(m{+}1)$-point
separation}
\begin{equation}
  \mathrm{sep}_{m+1}(X)\;=\;\max_{X'\subseteq X,\ |X'|=m+1}\ \ \min_{\bm{u}\neq\bm{v}\in X'}\|\bm{u}-\bm{v}\|_2,
  \label{eq:app-sep}
\end{equation}
with $\mathrm{sep}_{m+1}(X)=0$ when $|X|\le m$.

\begin{proposition}[Pigeonhole lower bound, every dimension]
\label{prop:pigeonhole}
For every $m\ge1$, every $d\ge1$ and every finite $X\subset\R^d$,
\begin{equation}
  q_m(X)\;\ge\;\tfrac12\,\mathrm{sep}_{m+1}(X).
  \label{eq:app-pigeonhole}
\end{equation}
\end{proposition}

\begin{proof}
If $|X|\le m$ both sides are $0$. Otherwise fix any $X'\subseteq X$ with
$|X'|=m+1$ and any codebook $C$ with $|C|\le m$. Assign each point of $X'$ to a
nearest centre; since $|X'|>|C|$, two distinct points $\bm{u},\bm{v}\in X'$ share a
centre $\bm{c}$. Then
$\|\bm{u}-\bm{v}\|_2\le\|\bm{u}-\bm{c}\|_2+\|\bm{c}-\bm{v}\|_2\le2\max_{\bm{x}\in X}\mathrm{dist}(\bm{x},C)$,
so $\max_{\bm{x}\in X}\mathrm{dist}(\bm{x},C)\ge\tfrac12\min_{\bm{u}\neq\bm{v}\in X'}\|\bm{u}-\bm{v}\|_2$.
Minimising the left side over $C$ and maximising the right over $X'$ gives
\eqref{eq:app-pigeonhole}.
\end{proof}

\begin{proposition}[Tightness at $d=1$, every $m$]
\label{prop:pigeonhole-tight}
For every $m\ge1$ and every finite $X\subset\R$, \eqref{eq:app-pigeonhole} holds
with equality: $q_m(X)=\tfrac12\,\mathrm{sep}_{m+1}(X)$.
\end{proposition}

\begin{proof}
Write $r=\tfrac12\mathrm{sep}_{m+1}(X)$ and $X=\{x_1\le\cdots\le x_T\}$. By
Proposition~\ref{prop:pigeonhole} it suffices to exhibit a codebook of at most $m$
centres with covering radius $\le r$. Construct one greedily: set $y_1=x_1$ and
$c_1=y_1+r$, which covers every point in $[y_1,y_1+2r]$; let $y_2$ be the smallest
point of $X$ exceeding $y_1+2r$, set $c_2=y_2+r$, and continue. The construction
terminates with some number $M$ of centres, and by definition it covers every point
of $X$ within distance $r$.

Suppose $M\ge m+1$. Then the first $m+1$ selected points $y_1,\dots,y_{m+1}$ satisfy
$y_{i+1}>y_i+2r$ for each $i$, so for $X'=\{y_1,\dots,y_{m+1}\}$ --- a subset of $X$
of size $m+1$, using $|X|\ge M\ge m+1$ --- every pair is separated by more than
$2r$, whence $\min_{\bm{u}\neq\bm{v}\in X'}|\bm{u}-\bm{v}|>2r$ and therefore
$\mathrm{sep}_{m+1}(X)>2r=\mathrm{sep}_{m+1}(X)$, a contradiction. So $M\le m$ and
$q_m(X)\le r$.
\end{proof}

\noindent Proposition~\ref{prop:pigeonhole-tight} is what licenses
Section~\ref{sec:theory}'s claim that the pigeonhole bound is not merely a bound but
\emph{is} $q_m$ for binary shared verbalizers: the budget bound then reads off the
geometry of the demands with no optimisation. We had previously only verified this
numerically ($600$ random instances, $n\in[2,7]$, $m\in[1,4]$, random scales, $0$
disagreements against the exact enumerator); the greedy argument above proves it.
The proof uses the ordering of $\R$ in an essential way --- the same greedy scheme in
$d\ge2$ is the standard factor-$2$ approximation and gives no equality --- and
consistently with that, at $d\ge2$ the bound is strict in the majority of random
instances ($164$ of $400$ strictly greater than the enumerated value). The
tightness claim must therefore stay confined to $d=1$;
Proposition~\ref{prop:pigeonhole} itself is valid in every dimension.

Both new statements above are checked numerically against an independently written
brute-force enumerator over contiguous blocks
(\texttt{experiments/phase2q\_bpo/tests/test\_appendix\_proofs.py}):
Proposition~\ref{prop:pigeonhole-tight} on $600$ random instances with
$n\in[2,8)$, $m\in[1,5)$ and scales spanning four orders of magnitude, exact to
$10^{-9}$; the greedy construction's centre count on a further $400$; and the
equilateral-triangle witness for strictness at $d=2$.

\subsubsection{Why the \texorpdfstring{$d=1$}{d=1} case is the one that matters here}

\begin{table}[t]
\centering
\small
\begin{tabular}{llcc c}
\toprule
label set & member tasks & $K_S$ & $d=K_S-1$ & scorable? \\
\midrule
\texttt{\{false, true\}} & WiC, QQP, BoolQA, MultiRC & $2$ & $\mathbf{1}$ & yes \\
\texttt{\{bad, good\}} & IMDB, SST-2 & $2$ & $\mathbf{1}$ & yes \\
\texttt{\{contradiction, entailment, neutral\}} & MNLI, CB & $3$ & $2$ & CB no \\
\texttt{\{negative, \dots, very positive\}} & Yelp, Amazon & $5$ & $4$ & neither \\
\bottomrule
\end{tabular}
\caption{Multi-member scopes on the $15$-task stream. Every scope that survives our
exclusion rules has $d=1$, so every $q_m$ we report is covered by
Propositions~\ref{prop:q2-closed} and~\ref{prop:pigeonhole-tight}. CB's rarest class
has $16$ examples, below the $64$/class the frozen protocol requires; Yelp and Amazon
fail the distinct-first-piece precondition, since \texttt{very positive} and
\texttt{very negative} share their first sentencepiece token.}
\label{tab:scope-dims}
\end{table}

Table~\ref{tab:scope-dims} is the reason the special case suffices, and the reason is
a property of the benchmark rather than a convenience of the analysis. It is also a
limitation in the direction we state in Section~\ref{sec:limitations}: the open case
is \emph{unmeasurable} here, not quietly assumed. A stream with a scorable $K_S\ge3$
shared scope would take the reported $q_m$ outside what we have proved.

\subsubsection{Transfer to accuracy, and the routing restriction}

\begin{corollary}[Budgeted accuracy bound, per-task routing]
\label{cor:qm-acc}
Under Assumption~\ref{asm:kappa-acc}, for any codebook of $m$ task-free offsets on
scope $S$ assigned \emph{per task} --- all examples of a task share a centre ---
\begin{equation}
  \max_{g\in\mathrm{mem}(S)}\big[R_g^{\mathrm{orc}}-R_g^{(m)}\big]
  \;\ge\;\kappa\cdot\min\big(q_m(Z_S),\,t_0\big),
  \label{eq:app-qm-acc}
\end{equation}
where $R_g^{(m)}$ is $g$'s accuracy under the centre it is assigned.
\end{corollary}

\begin{proof}
Identical to Proposition~\ref{prop:cheb-acc} with $r_S$ replaced by $q_m$: by
\eqref{eq:app-qm} some member is at distance at least $q_m$ from its assigned
centre, and \eqref{eq:app-cheb-acc-chain} applies verbatim to that member with
$t=\|\bm{c}(g)-P_0\bz_g\|_2\ge q_m$.
\end{proof}

\begin{remark}[The restriction is load-bearing: per-example routing breaks it]
\label{rem:per-example}
Corollary~\ref{cor:qm-acc} requires the router to be \emph{per task}. It is
\textbf{false} for a general per-example router, and not marginally so. Take one
binary task, $50$ examples with contrast $s=-3$ and true label $1$, and $50$ with
$s=+3$ and true label $0$. The best single offset scores $0.500$; two centres
$\{+4,-4\}$ with a router reading $\mathrm{sign}(s)$ score $1.000$ --- strictly
\emph{above} the single-offset oracle $R^{\mathrm{orc}}$, so no bound of the form
$R^{\mathrm{orc}}-R^{(m)}\ge(\text{positive})$ can hold. Per-example routing is a
strictly more powerful object than the one Proposition~\ref{prop:qm} bounds.

Our implementation contains both kinds. The prototype and batch-margin routers are
per task; the confidence router is per example, and it is the one measured to be
degenerate (oracle agreement $0.000$--$0.008$). So the theory and the reported
numbers are consistent --- but \eqref{eq:app-qm-acc} must be quoted with
``per-task routing'' attached, never as a statement about arbitrary task-free
routers. Whether a per-example task-free rule can be made to work is open, and
Remark~\ref{rem:per-example} shows it has \emph{more} headroom than the bound
allows, which makes the question live rather than closed.
\end{remark}

\subsection{Proof of Proposition~\ref{prop:amp} (margin-conditional amplification)}
\label{app:proof-amp}

This is the proposition behind Section~\ref{sec:theory}'s claim about rank budgets,
and it replaces a strict impossibility claim that an earlier version of this work
made and that is not correct. We state the repair first, then prove what survives.

\paragraph{What was claimed, and why it was wrong.}
The earlier assertion was that $\Delta^{\mathrm{id}}$ can be made large with
$\|\Delta\bW\|$ arbitrarily small, making rank-based forgetting bounds of the form
\eqref{eq:rankbound} vacuous for this term. That is false. A small $\|\Delta\bW\|$
\emph{does} imply a small logit perturbation, which \emph{does} imply a small
accuracy change --- unless many examples sit within that perturbation of the decision
boundary. The correct statement is conditional on the same margin density, and the
conclusion is weaker than impossibility.

\begin{proposition}[Margin-conditional amplification, restated]
\label{prop:amp-restated}
Let the adapter update change the hidden state by $\Delta\bh(\bx)=\Delta\bW\,\bh(\bx)$
for a matrix $\Delta\bW$ acting linearly on the state that feeds the unembedding.
Then for any two verbalizer tokens $\ell_j,\ell_k$ of a scope, the induced shift in
the pairwise logit difference satisfies
\begin{equation}
  \delta \;:=\; \sup_{\bx}\big|(\bu_{\ell_k}-\bu_{\ell_j})^\top\Delta\bh(\bx)\big|
  \;\le\; \|\bu_{\ell_k}-\bu_{\ell_j}\|_2\;\|\Delta\bW\|_2\;\sup_{\bx}\|\bh(\bx)\|_2 ,
  \label{eq:app-amp-delta}
\end{equation}
and, under Assumption~\ref{asm:kappa-dis} with $\delta\le t_0$, a shift of size
$\delta$ along the verbalizer-difference direction changes at least
$\kappa_{\mathrm{dis}}\,\delta$ of the predicted mass.
\end{proposition}

\begin{proof}
For \eqref{eq:app-amp-delta}, Cauchy--Schwarz and the definition of the spectral
norm give
$|(\bu_{\ell_k}-\bu_{\ell_j})^\top\Delta\bW\bh(\bx)|
\le\|\bu_{\ell_k}-\bu_{\ell_j}\|_2\|\Delta\bW\bh(\bx)\|_2
\le\|\bu_{\ell_k}-\bu_{\ell_j}\|_2\|\Delta\bW\|_2\|\bh(\bx)\|_2$;
take the supremum over $\bx$. For the second part, a shift of the pairwise logit
difference by $\delta$ is realised by displacing the gauge-fixed offset a distance
proportional to $\delta$ along the corresponding direction of $V_0$, so
Assumption~\ref{asm:kappa-dis} applies with $t=\delta$, giving flipped mass at
least $\kappa_{\mathrm{dis}}\min(\delta,t_0)=\kappa_{\mathrm{dis}}\delta$.
\end{proof}

\begin{remark}[Scope of \eqref{eq:app-amp-delta}, stated rather than glossed]
\label{rem:amp-scope}
Equation~\eqref{eq:app-amp-delta} is proved for an update entering \emph{linearly}
in the state that feeds the unembedding. For a LoRA update deeper in the network the
same inequality holds with an extra factor: if the map from adapter weights to the
final hidden state is $L$-Lipschitz in spectral norm on the relevant domain, then
$\delta\le\|\bu_{\ell_k}-\bu_{\ell_j}\|_2\,L\,\|\Delta\bW\|_2$. Our adapters are on
$q/v$ projections at every layer, so the deep case is the one that applies to our
runs and $L$ is not something we measure. This is where the unspecified constant $c$
of Section~\ref{sec:theory} lives, and we do not report a numerical value for it.
\end{remark}

\paragraph{The claim we are entitled to.}
Combining \eqref{eq:app-amp-delta} with the flip bound: a rank-$1$ update aligned
with a shared verbalizer-difference direction can satisfy the right-hand side of
\eqref{eq:rankbound} --- small $\sqrt{\rho}\|\bB\|\|\bA\|$ --- and still move
accuracy by order $\kappa\delta$. So \eqref{eq:rankbound} is loose by the factor
$\kappa$ in exactly the place that matters, and the entitled claim is neither
impossibility nor vacuity:

\begin{quote}
The rank bound is \textbf{uninformative about which of its two terms the damage
lands in}, and can be satisfied while $\Delta^{\mathrm{id}}$ is large.
\end{quote}

That is the operative point for allocators. Output-drift energy, importance scores
and rank minimisation are all averages over the whole drift; none separates the
component a zero-rank shared offset could absorb from the component that genuinely
requires subspace capacity. Under a fixed lifetime budget, rank spent on the former
is wasted --- and by Propositions~\ref{prop:cheb} and~\ref{prop:qm}, offset capacity
spent on conflicting tasks cannot substitute for rank either. Both budgets bind, for
different reasons.

\subsection{Claims we narrowed, and what forced each narrowing}
\label{app:narrowed}

Four statements in earlier versions of this analysis were stronger than what we can
prove. Each is recorded with the object that refuted it, because in every case the
correction changes how a number in the main text may be reported --- which is the
only reason to keep them on the page.

\paragraph{1. The candidate-centre enumerator is not exact beyond $d=1$.}
Our implementation computes $q_m$ by enumerating candidate centres over ``the points
together with all pairwise midpoints''. Step 2 of Proposition~\ref{prop:q2-closed}
shows this is exact at $d=1$, and also at $n\le2$ in any dimension, since the
midpoint of two points is their Chebyshev centre. It is \textbf{wrong for $d\ge2$}:
for three points forming a unit equilateral triangle the enumeration returns
$\sqrt3/2=0.8660$ while the true Chebyshev radius is the circumradius
$1/\sqrt3=0.5774$, an overestimate by exactly $1.5\times$, because the circumcentre
is neither a point nor a pairwise midpoint. No reported number is invalidated: the
enumerator is consulted only on scopes with $d=1$ or $n\le2$
(Table~\ref{tab:scope-dims}), and an overestimated radius inside a lower bound means
claiming less. But the docstrings and an earlier theory note asserted exactness
unconditionally, and that had to be narrowed to ``$d=1$ or $n\le2$''.

\paragraph{2. The pigeonhole bound is tight only at $d=1$.}
Proposition~\ref{prop:pigeonhole-tight} proves equality in one dimension for every
$m$. At $d\ge2$ the bound is strict in the majority of random instances --- $164$ of
$400$ --- so it must be quoted as a bound there, never as a formula.
Proposition~\ref{prop:pigeonhole} remains valid in all dimensions.

\paragraph{3. The accuracy-space bound needs a constant, and it is not the one we
measured.} Proposition~\ref{prop:no-free-lunch} shows no assumption-free
accuracy-space form exists. Beyond that, Lemma~\ref{lem:kappa-order} shows the
measured $0.0182$ is an upper estimate of the constant
Proposition~\ref{prop:cheb-acc} requires, because the estimator is a margin CDF and
therefore measures prediction disagreement rather than accuracy loss. The certified
statement is Lemma~\ref{lem:dis-minimax}. We rest no empirical claim on either.

\paragraph{4. Proposition~\ref{prop:qm} is proved for $d=1$, not in general.}
The case $m\ge2$ with $d\ge2$ is open. The correct phrasing of the gap is ``$m\ge2$
in dimension $\ge2$'', not ``$m\ge2$'': the $d=1$ case is closed for every $m$ by
Propositions~\ref{prop:q2-closed} and~\ref{prop:pigeonhole-tight}. We may write
``proved for binary shared verbalizers''; we may not write
``Proposition~\ref{prop:qm} is proved''.

\subsection{Status summary}
\label{app:status}

\begin{table}[h]
\centering
\small
\begin{tabular}{llll}
\toprule
statement & status & needs & where \\
\midrule
Prop.~\ref{prop:disjoint} (disjoint $\Rightarrow$ zero gap) & proved & --- & \S\ref{app:proof-disjoint} \\
Prop.~\ref{prop:cheb} (logit-space minimax) & proved & --- & \S\ref{app:proof-cheb} \\
Prop.~\ref{prop:q2-closed} (closed form $q_2$) & proved & $d=1$ & \S\ref{app:proof-qm} \\
Prop.~\ref{prop:pigeonhole} (pigeonhole bound) & proved & any $d$ & \S\ref{app:proof-qm} \\
Prop.~\ref{prop:pigeonhole-tight} (tightness) & proved & $d=1$, any $m$ & \S\ref{app:proof-qm} \\
Lem.~\ref{lem:dis-minimax} (disagreement minimax) & proved & Asm.~\ref{asm:kappa-dis} (measured) & \S\ref{app:proof-cheb-acc} \\
Prop.~\ref{prop:cheb-acc} (accuracy minimax) & proved & Asm.~\ref{asm:kappa-acc} (not measured) & \S\ref{app:proof-cheb-acc} \\
Cor.~\ref{cor:qm-acc} (budgeted accuracy) & proved & Asm.~\ref{asm:kappa-acc}, per-task routing & \S\ref{app:proof-qm} \\
Prop.~\ref{prop:amp} (amplification) & proved & linear case; $L$ unmeasured if deep & \S\ref{app:proof-amp} \\
$q_m$ for $m\ge2$, $d\ge2$ & \textbf{open} & --- & \S\ref{app:narrowed} \\
per-example task-free routing & \textbf{open} & --- & Rem.~\ref{rem:per-example} \\
\bottomrule
\end{tabular}
\caption{What is proved, under what, and where. Two rows are open and are reported as
open; the two conditional rows carry assumptions whose constants are discussed in
\S\ref{app:kappa-measured}, and only one of those constants is measured.}
\label{tab:proof-status}
\end{table}
